High-energy solar neutrinos from the pp-III chain are detected by Earth-based
experiments like Super-Kamiokande to study the Sun's core properties.

\begin{description}
\item[(a)] (5 points) The pp-III neutrino flux on Earth is
  $\Phi = 5.1 \times 10^6\,\mathrm{cm^{-2}\,s^{-1}}$. The Super-Kamiokande
  detector is a tank containing $5.0 \times 10^7\,\mathrm{kg}$ of pure water
  $\mathrm{H_2O}$. The molar mass of ($\mathrm{H_2O}$) is 18\,g/mol. Its
  detection efficiency is $\eta = 44.5\%$. Assume that solar neutrinos only
  interact with the detector through elastic scattering with electrons. The
  average neutrino-electron interaction cross-section is
  $\sigma \approx 1.0 \times 10^{-44}\,\mathrm{cm^2}$ in water. Physically,
  $\sigma$ can be understood as the effective target area a neutrino sees
  when approaching an electron; a larger $\sigma$ means a higher probability
  of collision. Calculate the number of events captured in a single day
  $N_\nu$ (24 hours).
\item[(b)] (6 points) Neutrinos ($\nu$) are detected indirectly via Cherenkov
  radiation emitted by scattered electrons ($\mathrm{e}^-$), i.e., a product
  of the electron-neutrino scattering $\nu + e^- \to e^- + \nu$ reaction, in
  water (refractive index $n = 1.33$). This event occurs when the velocity of
  electron satisfies $v \ge c/n$. Hence, we must use relativistic mechanics.
  The relativistic kinetic energy of electron is given as
  $K = (\gamma - 1)m_{\mathrm{e}}c^2$, where the Lorentz factor $\gamma$ is
  \[ \gamma = \frac{1}{\sqrt{1 - \frac{v^2}{c^2}}} \]
  The post-scattering energy of the neutrino ($E'_\nu$) is related to its
  pre-scattering counterpart ($E_\nu$) as
  \[ E'_\nu = \frac{E_\nu}{1 + \frac{E_\nu}{m_{\mathrm{e}}c^2}(1 - \cos\theta)} \]
  where $\theta$ is the scattering angle of the neutrino. Assuming the
  electron's initial kinetic and binding energies are negligible, determine
  the minimum initial neutrino energy ($E_{\nu,\mathrm{min}}$) in MeV required
  to generate a Cherenkov radiation event.
\item[(c)] (9 points) The solar pp-III reaction rate $R(T)$ depends on the
  core temperature $T$ via the Gamow equation
  $R(T) = R_0 e^{-BT^{-\frac{1}{3}}}$, where $R_0$ is a constant,
  $B = 17757\,\mathrm{K}^{1/3}$ is the Gamow factor, and
  $T = 1.5 \times 10^7\,\mathrm{K}$ is the steady-state core temperature.

  Suppose the temperature fluctuates slightly ($T \to T + \Delta T$), causing
  a small fractional reaction rate change ($R \to R + \Delta R$). The relation
  between the fractional change in the reaction rate and temperature is in
  the form $\frac{\Delta R}{R} = \alpha\frac{\Delta T}{T}$, where $\alpha$ is
  the thermal sensitivity coefficient. Express $\alpha$ in terms of $B$ and
  $T$. If the number of events detected by the Super-Kamiokande detector in a
  single day increases by 2\%, calculate the corresponding percentage change
  in the Sun's core temperature.

  \textit{Hints:} $(1 + x)^\delta \approx 1 + \delta x$ \textit{and}
  $\ln(1 + x) \approx x$ \textit{for} $|x| \ll 1$.
\end{description}
